% Physical page 5 onward: the developed commentary of the concise study,
% organized by the questions the three readings answer differently rather than
% by the order of the Mass. Each unit draws on several appointed elements and
% sets the readings' answers beside one another.

\section{The Propers: Detailed Commentary}
\zlabel{triptych:concise:commentary:start}

\subsection{What the father asked, and what he was given}

The Gospel turns on a rebuke. A man comes to ask his son's life
and hears, ``Unless you see signs and wonders, you believe not.'' St Gregory
the Great, preaching on this Gospel at the basilica of the martyrs Nereus and
Achilleus, says that the Gospel needs exhortation more than exposition and that only
this needs explaining. The man plainly believed, for \latin{neque enim ab eo
quaereret salutem quem non crederet salvatorem}: he would not have asked
salvation of one he did not believe a saviour. His request shows where the
faith failed. \latin{Poposcit namque ut descenderet, et sanaret filium ejus.
Corporalem ergo praesentiam Domini quaerebat, qui per spiritum nusquam deerat.}
He asked the Lord to come down, and so sought the bodily presence of one who,
\latin{per spiritum}, was nowhere absent; he honoured not majesty but bodily
presence. The Lord's answer corrects that faith by what it does: \latin{solo
jussu salutem reddidit qui voluntate omnia creavit}, he who created all things
by his will gave health by his command alone (\work{Homiliae in Evangelia} 28,
1). For Gregory the word is enough because of who speaks it.

St John Chrysostom grants the ruler a faith and finds it unsound: ``if this
ruler also believed, yet he believed not entirely or soundly, as is clear from
his enquiring `at what hour the fever left him'.'' He reads the rebuke as a
kindness aimed at the father: Christ ``toucheth his conscience, to show that His
miracles were wrought principally for the sake of the soul. For here He healeth
the father, sick in mind, no less than the son, in order to persuade us to give
heed to Him, not by reason of His miracles, but of His teaching''
(\work{Homilies on St John} 35, 2). Gregory teaches who it is that heals at a
distance; Chrysostom, what that healing does to the man who receives the word.
The first reading of the Mass is built on both, the object of faith and its
growth.

The Fathers do not agree on how much the man believed when he came. St
Augustine reads the rebuke at its full weight: Christ ``shows us a man
lukewarm, or cold in faith, or of no faith at all'', and since ``himself
believed'' comes only after the hour is known, ``it follows that when he was
making the request he did not yet believe'' (\work{Tractates on John} 16, 3).
St Cyril of Alexandria sees a faith real but childish. The ruler ``cometh as to
One able to heal, but he understandeth not yet that He is by Nature God''; yet
Christ ``doth not reject our lack of apprehension; but benefiteth even the
stumbling'', and ``The one command of the Saviour healeth two souls''
(\work{Commentary on John} II.5). Bl.\ Ildefonso Schuster, expounding this
Mass, says that the ruler ``had faith in Jesus'', which Christ desired ``first
of all to purify \dots\ from any material element and from all idea of
self-interest'', requiring him ``to believe in his word without actually
beholding the cure'' (\work{The Sacramentary} III, p.~176). St Thomas Aquinas
gathers the answers into a sequence. The man's faith was weak when he first
interceded, firmer when he called Jesus Lord, more perfect when he believed the
word and went, and perfected when he learned the hour and believed with his
house; the evangelist's \latin{et ibat} means that he went on growing, because
on the way of God not to go forward is to go back (\work{Super Ioannem} c.~4,
lect.~7). So Augustine and Gregory differ on the faith that came, and
Chrysostom, Cyril, Schuster and Aquinas stand between them; but Augustine,
Chrysostom, Cyril and Aquinas agree that a fuller faith followed the word, and
the Gospel's last clause bears them out.

A second question sharpens the first. The same Lord who would not go down to
the ruler's son offered to go to the centurion's slave (Mt 8:7). Gregory reads
the contrast as a rebuke to pride, which honours in men their wealth and rank
and not the nature made to God's image: \latin{Ecce de coelo venit qui servo in
terra occurrere non despicit}, he who came from heaven does not disdain to go
to meet a slave on earth (28, 2). Chrysostom answers from the two men's faith:
``in the former case faith had been perfected, and therefore He undertook to
go'', while the ruler ``knew not yet clearly that even when absent He could
heal''. Augustine answers from their humility, and so opens the third reading:
the centurion ``declared himself unworthy of His presence'', the ruler ``sought
His presence by force''; ``Here is a ceding to loftiness; there, a conceding to
humility'' (16, 5). For Gregory the difference lies in the honour men give, for
Chrysostom in a faith not yet complete, for Augustine in the lowliness that
believes a word can work at a distance.

The rebuke is hard, and the witnesses weigh it differently. Chrysostom, Cyril
and Aquinas hold that it served the man; Aquinas adds,
after 1 Cor 14:22, that it fell on one instructed in the law who sought to
believe by signs, since signs are given to unbelievers and the faithful are led
by the Scriptures. Augustine does not soften it, and sets the man beside St
Thomas the Apostle, to whom it was said, ``Blessed are they that have not seen,
and yet have believed'' (16, 4). Signs are not made worthless. The healing is itself a sign, the second at
Cana, and Gregory says elsewhere that visible miracles were given \latin{ut
corda videntium ad fidem invisibilium pertrahant}, to draw the hearts of those
who saw them to faith in what is unseen (\work{Homiliae in Evangelia} 4, 3).
The father was given both, in an order that matters: the word first, believed;
the sign after, found true at the hour of the word. ``What now are we taught by
these things?'' Chrysostom asks. ``Not to wait for miracles, nor to seek pledges
of the Power of God'' (35, 3). He draws that lesson into a household's life:
many give thanks when a child recovers, and they ought to give thanks ``even if
they obtain it not''.

\subsection{One household, or two peoples}

Augustine reads the healing a second way, in its chapter, and the third reading
is built on that sequence. His sixteenth \work{Tractate} begins from a puzzle:
going from Samaria, where he was honoured, into Galilee, where he was brought
up, Jesus ``gave testimony that a prophet hath no honour in his own country''
(Jn 4:44). ``In Samaria He spent two days, and the Samaritans believed on Him;
many were the days He spent in Galilee, and yet the Galileans did not believe on
Him.'' ``At His discourse alone many of the Samaritans believed; at that
miracle, in the place where it was wrought, only that house believed'' (16, 3).
A \latin{prodigium}, he says, is so called \latin{quod porro dicat}, because it
betokens something to come, and these events betoken the order of salvation:
``Let us, I say, assume the Lord's own country to mean the people of the Jews.''
Then he turns to his hearers: ``I am speaking to God's people; so many of us
have believed, what signs have we seen?'' The blessing spoken to St Thomas goes
to them: ``We are the persons here foretold'' (16, 4).

Aquinas sets the father on the other side of the same history. Reading the
episode allegorically, he makes the ruler Abraham, or one of the Fathers of the
Old Testament, and the son the people of Israel, sick with depraved pleasures
and teachings in Capharnaum, whose name he renders ``abundance'': the abundance
of which Moses sang, ``The beloved grew fat, and kicked'' (Deut 32:15). The
ruler's plea is the Fathers' intercession for their people, as Jeremias is
called ``a lover of his brethren, and of the people of Israel'' (2 Mac 15:14);
the cure comes by the word at the seventh hour, which signifies the seven gifts
of the Holy Spirit, by whom sins are forgiven; and the slowness of the
Galileans beside the Samaritans is the slowness with which Israel came to the
faith, little by little, like grapes gleaned after the vintage (Mic 7:1).
Augustine and Aquinas agree that the episode figures the relation of Israel and
the nations to Christ's word and that Israel's patriarchs keep their honour.
They disagree on what the father figures: for Augustine he belongs to the
country that needed signs, for Aquinas he is the patriarch who intercedes for
it. St Cyril gives a third identification. The evangelist's aim, he says, was
``to shew how superior in obedience were the aliens to the Jews'', and the
ruler, who ``believed with his whole house'', stands beside the Samaritans
among the obedient.

Aquinas himself shows how the first and third readings relate. The same
lecture that traces one man's faith from weakness to its perfection reads him,
at another sense of the text, as the Fathers praying for their people. The two
readings answer different questions of the same words: what faith the word
asks of a believer, and whose history the healing figures. Neither denies the
other.

The third reading can be heard as setting Israel against the Church, and its
witnesses say both more and less than that. Augustine draws St Paul's olive
from the centurion: the proud branches broken, the humble wild olive grafted
in, ``nevertheless let the root remain \dots\ Where does the root remain? In
the patriarchs''; and they are ``In rest with God, in great honor'' (16, 5). He
ends with the country Christ founded, from Sion, of the Virgin Mary: ``He had
no honor in His country, wherein He was formed; let Him have honor in the
country which He has formed'' (16, 7). Cyril, who says that Israel ``with
reason falleth from the hope'', ends in hope: those set in the first rank ``because of the election of the
fathers, will come last and after the calling of the Gentiles. For when the
fulness of the Gentiles is come in, then shall all Israel be saved''
(\work{Commentary on John} II.5, after Rom 11:25--26). Augustine's conclusion is humility for the grafted branch,
which stands by faith and not by birth; Aquinas's, the intercession of the
Fathers for their people; Cyril's, hope.

The Gospel's last words, \latin{et credidit ipse, et domus eius tota}, show
where the readings meet and part. In the first they end one man's growth in
faith, from need, through trust in a word, to belief with his house. In the
third they figure the whole house: the people whose Fathers prayed for it and
the nations grafted beside it. The second reading does not take up those words.
It reads the father's journey as one exile's crisis in a world that passes.
Gregory preached in a time when, he says, \latin{Ubique mors, ubique luctus,
ubique desolatio}, and the world that once held men by delight \latin{tantis
plagis plenus est, ut ipse nos jam mundus mittat ad Deum}, is so full of blows
that it sends us to God itself (28, 3). That world sends the father to Christ, and the word that reaches his son is
the word St Jerome finds in the furnace: the angel who goes down into the fire
is \latin{sermo divinus}, the divine word, which comes to the soul that has
despaired of human help and turned wholly to the Lord (\work{Commentarii in
Danielem}, on 3:49).

\subsection{The confession at the entrance}

All three readings begin at the Introit, and each hears its confession
differently. Two of them hear it through St Augustine's account of how the
saints pray. Expounding the Gradual's psalm at ``The Lord is just in all his
ways'', he writes: \latin{omnes sancti in tribulationibus constituti, prius
justitiam ejus laudaverunt, et sic beneficia petiverunt. Prius dixerunt: Justum
est quod facis. Sic rogavit Daniel}: all the saints, set in tribulation, first
praised his justice and so asked his benefits; so Daniel asked
(\work{Enarrationes in Psalmos} 144, on v.~17). ``First they praised Him
scourging,'' runs the English, ``and so they felt Him feeding.'' The first
reading takes that order as the beginning the father's faith needed: with God
and his justice, not with the sign. The second takes it as the exile's first
word, the confession of a people that has lost its city.

St Jerome gives that confession its rule and its limit. On Azarias's first
verse he writes, \latin{Quando diversis premimur angustiis, ex toto cordis hoc
loquamur affectu, et quidquid nobis acciderit, juste nos sustinere fateamur}:
when we are pressed by distresses of every kind, let us confess with the whole
heart that whatever has befallen us we bear justly. But \latin{et certe tres
pueri non peccaverant}: the three youths had not sinned, and were not of an age
to be punished for their own faults when they were carried to Babylon;
\latin{ergo \dots\ ex persona populi loquuntur}, they speak in the person of the
people (\work{Commentarii in Danielem}, on 3:26, 29). Jerome closes the hard
reading of the Introit, that all suffering is the sufferer's own punishment;
Augustine sets the confession in the order of prayer, not in an account of the
causes of every pain. And Jerome gives the Church leave to pray the prayer as
hers, \latin{si quando Ecclesiae propter peccata populi \dots\ sustinent
penuriam, et quando in persecutionibus non offertur sacrificium et oblatio},
whenever the Churches suffer want for the people's sins, and when in
persecution the sacrifice is not offered (on 3:37).

Here the second and third readings meet. For the second, the innocent youths
teach the Church in exile to confess God's justice with and for a people. For
the third, the Introit is Israel's own confession, which the Church prays as
hers in the way Azarias prayed it, for a people whose sin he had not committed.
Jerome's \latin{ex persona populi} serves both.

The Introit's verse, \latin{Beati immaculati in via: qui ambulant in lege
Domini}, is heard three ways. For the first reading the way is Christ. St
Hilary of Poitiers recalls \latin{Ego sum via} and insists that the way is
\latin{non solum ineunda \dots\ sed etiam pergenda}, to be walked to its end and
not only entered (\work{Tractatus super Psalmos}, in Ps.~118, \latin{Aleph} 2),
and Cassiodorus takes the way to be the Lord Saviour, walked with the steps of
good action. Neither speaks of the Gospel; but the father walks such a road,
\latin{Credidit homo sermoni \dots\ et ibat}, before he knows what he will find
at its end. For the second reading the verse names the exiles' road home.
Augustine makes it the divine voice recalling everyone who seeks blessedness
by evil roads: \latin{si ergo vis esse beatus, esto immaculatus. \dots\ Sed ubi
erit quisque immaculatus, nisi in via? In qua via, nisi in lege Domini?} and to
the lustful, the avaricious and the ambitious, \latin{Quo itis? Peritis, et
nescitis} (\work{Enarrationes in Psalmos} 118, sermo 1, 1). Hilary reads the
verse as the first stage of an ascent, in which the undefiled curb the flesh
and conquer avarice before they search God's testimonies; St Robert Bellarmine
sets the law of the Lord against the way of concupiscence, ``full of the dust of
pride, the mud of luxury, and the dirty water of avarice''. The road out of
Babylon is the keeping of the commandments the exiles confess they broke. For
the third reading Cassiodorus supplies another voice: the whole of Ps~118 is
spoken by the single chorus of the saints, from the world's beginning to its
end, set forth as one speaker for the unity of the Church. The blessing of
those who walk in the law is then spoken by the saints of both peoples
together.

The orations carry the confession forward. Schuster expounds the order of the
Collect's two gifts: ``Peace comes after mercy, for as long as grace has not
blotted out the sin, the heart, torn by remorse, weakened by its own passions
and at war with itself, cannot find peace'' (p.~175). The first reading hears
in \latin{secura mente} the mind the father went home with, at rest before the
proof arrived. The Collect asks that mind through pardon, and Cassiodorus says
that the faithful say the Alleluia's \latin{Paratum cor meum} rightly when they
trust that their sins are forgiven by God's gift. The
second hears the orations of exiles who have confessed: Sicard of Cremona and
William Durandus, expounding these chants of captivity, say that the bodily
captivity signifies our spiritual one and that the return from it is the
remission of sins (\work{Mitrale} VIII.20; \work{Rationale} VI.137). The third
hears \latin{fidelibus tuis} as the faithful of every people grafted in. And
the Postcommunion's \latin{oboedire mandatis}, the one answer all three
readings give the Introit's \latin{non oboedivimus}, is in the first the word
received becoming the word obeyed, in the second the exiles' road home, and in
the third the law given to Israel asked now of all who receive the gifts. Dom
Lucien Fromage, in the continuation of \work{The Liturgical Year}, joins it to
the altar: ``A persevering fidelity in observing God's commandments is the best
preparation a christian can make for approaching to the holy Table.''

\subsection{By the rivers of Babylon}

St Augustine preached on the Offertory's psalm on a day when it had been sung,
and the second reading takes its programme from that sermon. ``The waters of
Babylon are all things which here are loved, and pass away.'' One man loves
farming, another soldiering, another office; let each see ``that what he hath
loved is not a foundation of Jerusalem, but a stream of Babylon''. The citizens
of Jerusalem do not throw themselves in: ``Sit `by' the waters, not `in' the
waters, not `under' the waters; but yet sit, in humble fashion.'' Their tears
are not Babylon's: ``If thou weepest in the remembrance of Sion, thou oughtest
to weep even when it is well with thee in Babylon.'' Their instruments are
``God's Scriptures, God's commands, God's promises, meditation on the life to
come'', hung on the willows, barren trees watered by Babylon, because such men
are unfit to carry them (\work{Enarrationes in Psalmos} 136, on vv.~1--2). He asks that the
song be done as well as sung, \latin{oportet ut non tantum ista decantemus, sed
et faciamus}, and ends, \latin{Fratres, organa in operando non cessent; vobis
invicem cantate cantica Sion}: brothers, let your instruments not rest in your
work; sing to one another the songs of Sion.

St Hilary reads the same psalm more inwardly. Its captivity is spiritual; Sion
is the seat of eternal blessedness, \latin{mater coelestis nobis Jerusalem},
and man has been exiled from her since Adam's sin. The rivers are the works of
the age and of the body, which flow past; the willows, which revive in water,
are the saints, revived by the word of God and the water of baptism (Isa 44:4);
the instruments are our bodies, hung up in continence (\work{Tractatus super
Psalmos}, in Ps.~136). Augustine and Hilary disagree within a shared reading.
For Augustine, and for Bellarmine after him, the willows are barren men; for
Hilary, and for Cassiodorus, they are saints. The instruments are the
Scriptures for Augustine and our bodies for Hilary, and Cassiodorus takes
Augustine's instruments and Hilary's willows, the Scriptures hung on holy men.
Cassiodorus draws a posture from \latin{super flumina}, the freed sitting on the
banks and not in the stream, and calls the singing of psalms the consolation of
those who weep and a remedy of the soul.

The other two readings hear the chant differently. For the first, the Offertory
weeps for a Sion remembered and not seen, and faith of the father's kind lives
on a word held in absence. For the third, the chant is first Israel's.
Fromage gives the Mass's chants to two choirs: ``The Gentile-Choir takes the
Gradual and Communion-Anthem; the Choir of the Jews, the Introit and
Offertory''. ``In the Introit, the Jewish people sang its repentance and humble
confidence'', and in the Offertory Israel, grown repentant, ``joins our Mother
the Church'' in the psalm of exile (\work{The Liturgical Year} XI, pp.~439,
447, 451). Here the readings meet a real tension. Fromage hears Israel's voice
first in the Offertory; Augustine's sermon hears in the same verse the Church's
own captivity. Both are traditional answers, and they differ on whose exile
the chant first voices. Jerome's youths, who confess for a people they belong
to, show how one voice can carry both.

The second reading alone faces the end of the psalm, which the Mass does not
sing: ``Blessed be he that shall take and dash thy little ones against the
rock'' (Ps 136:9). Augustine asks, ``What are the little ones of Babylon? Evil
desires at their birth'', and answers the fear that the dashed may not die:
``Dash it against the Rock; and that Rock is Christ.'' Hilary reads the little
ones as \latin{tenera adhuc corporis vitia}, the body's vices while still
tender, since \latin{Petra autem, secundum Apostolum, Christus est}
(in Ps.~136, 14), and Cassiodorus as the flesh's errors while they are small.
Bellarmine makes the little ones persons: Babylon's young, converted and
brought ``to the rock, Christ, to get a fortunate dash against it, and die the
death of the old man, to rise a new man''; and, morally, those not yet deeply
stained, for ``blessed is he, who, by a happy dash on the rock, kills sin''.
Chrysostom reads no figure in the verse at all: the words are the captives' own
anger, which the prophet reports and does not make his own, and the New
Covenant commands that we feed our enemies and pray for those who abuse us.
Four witnesses make the rock Christ; three of them make the little ones sins
at their birth, and the fourth persons brought to Christ to die to sin.
Chrysostom alone leaves the verse as the captives' passion and answers it with
the Gospel.

\subsection{How to sing the Lord's song in a strange land}

The captives ask, \latin{Quomodo cantabimus canticum Domini in terra aliena?}
The Epistle answers them before the Offertory is sung, and the Alleluia
answers in the Epistle's own verbs, \latin{cantantes, et psallentes} there,
\latin{cantabo, et psallam} here. St Jerome reads ``redeeming the time'' as
the wise man buying back what malice had sold, \latin{dies malos in bonos
vertimus}, we turn evil days into good, as no change of times changed Joseph
or Job. ``Singing and making melody in your hearts'' is a warning to singers:
\latin{Deo non voce, sed corde cantandum}, God is to be sung to not with the
voice but with the heart, and a singer of poor voice whose works are good
\latin{dulcis apud Deum cantor est}. ``Giving thanks always for all things''
reaches loss: if a house falls or wife and children are taken captive, the
Christian says, \latin{Benedictus Deus, minora me scio sustinere quam mereor},
blessed be God, I know that I suffer less than I deserve (\work{Commentarii in
Epistolam ad Ephesios} III, on 5:16, 19, 20). Chrysostom gives the same verses
the exile's turn: ``The time is not yours. At present ye are strangers, and
sojourners, and foreigners, and aliens''; ``only preserve the principal thing,
I mean the faith''. Against drunkenness he sets the psalms, ``For they who sing
psalms are filled with the Holy Spirit'', and he reads ``with your hearts'' as
``with close attention and understanding'' (\work{Homilies on Ephesians} 19).
Jerome subordinates the voice to the heart and Chrysostom asks for the voice
with attention; both oppose singing that seeks only to please, and Bellarmine,
on the Offertory's psalm, makes that the fault of those who sing the Lord's
song in a strange land.

The first reading hears a different demand in the Epistle. Its
\latin{intellegentes quae sit voluntas Dei} is, for Aquinas, the first
principle of the prudence the Apostle asks, as the Lord's Prayer asks
\latin{Fiat voluntas tua} (\work{In Epistolam ad Ephesios} c.~5, lect.~6). The
father's going was such prudence: he acted on a will declared to him in a
word. The third reading hears the Epistle's addressees: Gentile Ephesians who
were ``heretofore darkness, but now light in the Lord'' (Eph 5:8), and whose new
life the lesson describes.

The Alleluia's verse is also Ps~56:8, and one sentence of Augustine on it
serves two readings in two ways: ``the patience of good men with preparation of
heart accepteth the will of God''. The first reading takes the acceptance.
Augustine sets the prepared heart against what others prepare, ``He hath
prepared pit to deceive, shall I not prepare heart to suffer?'', and Bellarmine
makes it a whole surrender, ``ready to take anything cheerfully from your
hand''. The father had fixed the means of his son's healing, that Christ should
come down, and was given another, a word in place of a presence; he took it and
went, with the heart the Alleluia sings, as Chrysostom asks every household to
give thanks ``even if they obtain it not''. The second reading takes the song.
Augustine goes on, ``and glorieth in tribulations'', and his example is St Paul
in chains: ``Whence, but that prepared was his heart? Therefore he was singing
and playing'' (\work{Enarrationes in Psalmos} 56, on v.~8). The captives could
not sing the Lord's song for their captors; the prepared heart sings it to God,
\latin{psallam tibi}. Schuster reads the Missal's \latin{gloria mea} of God, the
soul's glory ``also because he alone is the just judge of our merits''. The
glory the verse sings is God himself, and not the honour that, Gregory says,
pride pays to rank and Christ would not pay to the ruler's.

The third reading follows the psalm one verse further than the Mass: ``I will
praise thee, O Lord, among the people: and I will sing unto thee among the
nations'' (Ps 107:4). Cassiodorus reads that confession among the peoples and
that song among the nations as the Church gathered first from the Jews and
then from the Gentiles.

\subsection{Fed in due season, and a word remembered}

The Gradual and the Communion share a posture: those who pray have been
promised something and have not yet received it. Augustine reads the
Gradual's due time as a physician's care, ``Just as when thou refreshest a
sick man in due season, when he ought to receive, then Thou givest'', and says
to anyone who has asked and not received, ``let his eyes wait for the food,
which He giveth in due season \dots\ Thou delayest, not deniest''
(\work{Enarrationes in Psalmos} 144, on vv.~15--16). The first and second
readings both use the sentence, for different waitings. In the first it is the
father's night between the word and the news, \latin{heri hora septima}, the
food come in its season. For that faith in a word unseen, Schuster reads the
Gradual's food as the Eucharist: the bread
God prepares ``in every region of the earth'' is given under ``the veil of faith,
in order that we may have all the merit of believing in the pure and single
word of the Son of God'' (pp.~175--176). In the second it is the exiles', fed
and waiting. For the third reading the Gradual's words are universal,
\latin{Oculi omnium} and \latin{omne animal}, and Fromage gives the chant to the
Gentile choir.

The Communion asks God to remember his word, and its witnesses take the
request as a prayer that God fulfil a promise. Augustine: ``When
therefore it is said unto Him, `O remember,' the desire of him who prayeth is
displayed, because he asketh for what was promised \dots\ that is, fulfil Thy
promise to Thy servant'' (118, sermo 15, on v.~49). Bellarmine adds why the
promise must be asked: God ``still wishes his faithful servants to ask him to
carry them out; and thus, prayer becomes one of the means through which God
decreed to fulfill his promises''. The readings then hear different speakers.
In the first it is the father, whose double petition was answered by a word to
be believed. In the second it is the Church in exile. Augustine reads the
humiliation as tribulation and says, ``This hope was given to Christ's Body,
that is, to the Church, that it might be a comfort to Her in her humiliation''
(on v.~50), and Hilary hears in it the comfort every word of God gives in
sickness, persecution, loss and bereavement, since each calls us to the hope of
heavenly goods. In the third it is the Church of the nations, hoping in a word
first given to Israel's father. Bellarmine names the promise: ``when you said
to Abraham, and through him to all his children, `Walk before me, and be
perfect;' I will be `thy reward exceedingly great'.'' The verse the Communion
leaves unsung gives the first reading its last link: \latin{quia eloquium tuum
vivificavit me}, because thy word has given me life, beside the Gospel's
\latin{Filius tuus vivit}.

The Secret carries the father's situation to the altar. Fromage joins it to
this Gospel: ``The whole power of the God, who, with a word, cures both soul
and body, resides in the Mysteries which are about to be celebrated on our
Altar here''; and, of the father's plea that Christ come down, ``our Jesus has no need to
come down from heaven to earth, in order to give efficiency to the commands of
his gracious will'' (\work{The Liturgical Year} XI, p.~451). Schuster asks of
the same prayer ``that the holy Sacrament may be to us a spiritual medicine and
antidote, to protect us from the virus of sin which poisons our blood''
(p.~177). The second reading sets the Secret's medicine beside Cassiodorus's
remedy of the soul in the exiles' psalmody; the third, beside the cure that
Aquinas's sick son, the people grown fat on abundance, needed.

The readings end in three hopes, and each is its own. For the first, the
father's road home between the word and the sight of its fulfilment is the
life of faith, which ends when the promise is kept and seen; Schuster says that
``the blessed enjoy that same truth which sustains us here on earth; but they
enjoy it without any veil, face to face'' (p.~176). For the second, the memory
of Sion is longing for the city: ``whither your hope goeth before, let your
life follow \dots\ Your captivity will pass away, your happiness will come; the
last enemy shall be destroyed, and we shall triumph with our King, without
death'' (Augustine, on Ps~136). For the third, the healing of a whole house at
Capharnaum looks to the healing of both peoples in the one house of God:
``then shall all Israel be saved'' (Cyril). The second reading makes no claim
about the other two, but set beside them it gives them a setting. The father
who walks home on a word, and the nations who believe without signs, are both
on the road of exiles who remember Sion.
